%%%%%%%%%%%%%%%%%%%%%%%%%%%%%%
% Biblioteca de Glossários,
% lista de siglas e símbolos
% preencha aqui com suas siglas e símbolos
% Neste Template o Glossário será apenas
% para compilações de PFC e Pré-Projetos
%%%%%%%%%%%%%%%%%%%%%%%%%%%%%%
%
%%%%%%%%%%%%%%%%%%%%%%%%%%%%%
%   SIGLAS
%%%%%%%%%%%%%%%%%%%%%%%%%%%%%
%
\newglossaryentry{IEEE} % 
{
  name = {IEEE},
  description = {IEEE - Instituto de 
      Engenheiros Eletricistas e Eletrônicos (\textit{IEEE is the world's largest technical professional organization dedicated to advancing technology for the benefit of humanity})}
}
%%%
\newglossaryentry{ONS}
{
  name = {ONS},
  description = {Operador Nacional do 
      Sistema}
}
%%%%%
\newglossaryentry{CAPES}
{
   name = {CAPES},
   description = {Centro de Aperfeiçoamento 
       de Pessoal de Nível Superior}
}
%%%%%
\newglossaryentry{MPPT}
{
   name = {MPPT},
   description = {\textit{Maximum Power 
       Point Tracking}}
}
%%%%%
\newglossaryentry{COEEL}
{
   name={COEEL},
   description={\textit{Coordenação do Curso de Engenharia Elétrica do IFBA campus Vitória da Conquista}}
}
%%%%%
\newglossaryentry{ISO}
{
   name={ISO},
   description={\textit{International Organization for Standardization (Organização Internacional para Padronização)}}
}
%%%%%
\newglossaryentry{CIE}
{
   name={CIE},
   description={\textit{Commission Internationale de l'éclairage (Comissão Internacional de Iluminação)}}
}
%%%%%
\newglossaryentry{ABNT}
{
   name={ABNT},
   description={\textit{Associação Brasileira de Normas Técnicas}}
}
%%%%%
\newglossaryentry{NBR}
{
   name={NBR},
   description={\textit{Norma Brasileira}}
}
%%%%%
\newglossaryentry{NHO}
{
   name={NHO},
   description={\textit{Norma de Higiene Ocupacional}}
}
%%%%%
\newglossaryentry{PID}
{
   name={PID},
   description={\textit{Proporcional, Integrativo e Derivativo}}
}
%%%%%
\newglossaryentry{HTTP}
{
   name={HTTP},
   description={\textit{Hypertext Transfer Protocol (Protocolo de Transferência de Hipertexto)}}
}

%%%%%
\newglossaryentry{EPE}
{
   name={EPE},
   description={\textit{Empresa de Pesquisa Energética}}
}

%%%%%
\newglossaryentry{DSP}
{
   name={DSP},
   description={\textit{Digital Signal Processors (Processadores de Sinais Digitais)}}
}

%%%%%
\newglossaryentry{FPGA}
{
   name={FPGA},
   description={\textit{Field-Programmable Gate Arrays (Matrizes de Portas Programáveis)}}
}

\newglossaryentry{USB}
{
   name={USB},
   description={\textit{Universal Serial Bus (Barramento Serial Universal)}}
}

\newglossaryentry{UART}
{
   name={UART},
   description={\textit{Universal Asynchronous Receiver/Transmitter (Transmissor/Receptor Assíncrono Universal)}}
}

\newglossaryentry{USART}
{
   name={USART},
   description={\textit{Universal Synchronous/Asynchronous Receiver/Transmitter (Transmissor/Receptor Síncrono/Assíncrono Universal)}}
}

\newglossaryentry{I2C}
{
   name={I2C},
   description={\textit{Inter-Integrated Circuit (Circuito Inter-Integrado)}}
}

\newglossaryentry{SPI}
{
   name={SPI},
   description={\textit{Serial Peripheral Interface (Interface Periférica Serial)}}
}

\newglossaryentry{CAN}
{
   name={CAN},
   description={\textit{Controller Area Network (Rede de Área do Controlador)}}
}
%%%%%
\newglossaryentry{PWM}
{
   name={PWM},
   description={\textit{Pulse Width Modulation (Modulação por Largura de Pulso)}}
}
%%%%%
\newglossaryentry{TRIAC}
{
   name={TRIAC},
   description={\textit{Triodo de Corrente Alternada}}
}
%%%%%
\newglossaryentry{CA}
{
   name={CA},
   description={\textit{Corrente Alternada}}
}
%%%%%
\newglossaryentry{CC}
{
   name={CC},
   description={\textit{Corrente Contínua}}
}
%%%%%
\newglossaryentry{MT1}
{
   name={MT1},
   description={\textit{Main Terminal 1 (Terminal Principal 1)}}
}
%%%%%
\newglossaryentry{MT2}
{
   name={MT2},
   description={\textit{Main Terminal 2 (Terminal Principal 2)}}
}
%%%%%
\newglossaryentry{G}
{
   name={G},
   description={\textit{Gate (Porta)}}
}
%%%%%
\newglossaryentry{PyQt5}
{
name={PyQt5},
description={\textit{Conjunto de bindings do toolkit Qt~5 para Python.}}
}
%%%%%
\newglossaryentry{BLE}
{
name={BLE},
description={\textit{Bluetooth Low Energy (Bluetooth de Baixo Consumo).}}
}

%
%%%%%%%%%%%%%%%%%%%%%%%%%%%%
% SÍMBOLOS
%%%%%%%%%%%%%%%%%%%%%%%%%%%%
%

\newglossaryentry{FP}
{
   name={$cos(\phi)$},
   description={Fator de Potência}
}
%%%%%
\newglossaryentry{fi}
{
   name={$\phi$},
   description={Ângulo do Fator de Potência}
}
%%%%%
\newglossaryentry{N1}
{
   name={$N_1$},
   description= Número de espiras do primário
}
%%%%%
\newglossaryentry{U1}
{
   name={$U_1$},
   description= Tensão do primário (V)
}
%%%%%
\newglossaryentry{freq}
{
   name={$f$},
   description= frequência (Hz)
}
%%%%%
\newglossaryentry{SL}
{
   name={$S_L$},
   description= Área da Seção Líquida do núcleo ($cm^2$)
}
%%%%%
\newglossaryentry{DenFlux}
{
   name={$B$},
   description= Densidade de Fluxo Magnético (Gauss)
}
%%%%%
\newglossaryentry{Ua}
{
   name={$U_a$},
   description= Tensão na Fase A
}
%%%%%
\newglossaryentry{Ia}
{
   name={$I_a$},
   description= Corrente na Fase A
}
%%%%%
\newglossaryentry{w}
{
   name={$\omega$},
   description= Velocidade angular $(rad/s)$ 
}
%%%%%
\newglossaryentry{lux}
{
   name={$lx$},
   description=  Intensidade luminosa
}
%%%%%
\newglossaryentry{kelvin}
{
   name={$K$},
   description= Temperatura de cor
}
%%%%%
\newglossaryentry{UGR}
{
   name={$UGR$},
   description={\textit{Índice de Ofuscamento Unificado}}
}
%%%%%
\newglossaryentry{UGRL}
{
   name={$UGRL$},
   description={\textit{Índice de Ofuscamento Unificado Limite}}
}
%%%%%
\newglossaryentry{Ra}
{
   name={$R_a$},
   description={\textit{Índice geral de reprodução de cor}}
}
%%%%%
\newglossaryentry{Tcp}
{
   name={$T_{cp}$},
   description={\textit{Temperatura de cor correlata}}
}
%%%%%
\newglossaryentry{Kp}
{
   name={$K_{p}$},
   description={\textit{Ganho proporcional}}
}
%%%%%
\newglossaryentry{Ki}
{
   name={$K_{i}$},
   description={\textit{Ganho integral}}
}
%%%%%
\newglossaryentry{Kd}
{
   name={$K_{d}$},
   description={\textit{Ganho derivativo}}
}
%%%%%
\newglossaryentry{uk}
{
   name={$u(k)$},
   description={\textit{Sinal de controle no instante $k$}}
}
%%%%%
\newglossaryentry{ek}
{
   name={$e(k)$},
   description={\textit{Erro no instante $k$}}
}
%%%%%
\newglossaryentry{Ts}
{
   name={$T_{s}$},
   description={\textit{Tempo de amostragem}}
}
%%%%%
\newglossaryentry{Ls}
{
   name={$L_{s}$},
   description={\textit{Atraso do sistema}}
}
%%%%%
\newglossaryentry{tau}
{
   name={$\tau$},
   description={\textit{Constante de tempo do sistema}}
}
%%%%%
\newglossaryentry{Ku}
{
   name={$K_{u}$},
   description={\textit{Ganho Crítico}}
}
%%%%%
\newglossaryentry{Tu}
{
   name={$T_{u}$},
   description={\textit{Período de oscilação sustentada}}
}
%%%%%
\newglossaryentry{TON}
{
   name={$T_{ON}$},
   description={\textit{Tempo em que o sinal permanece em nível alto (ligado)}}
}
%%%%%
\newglossaryentry{TOFF}
{
   name={$T_{OFF}$},
   description={\textit{Tempo em que o sinal permanece em nível baixo (desligado)}}
}
%%%%%
\newglossaryentry{P}
{
   name={$P$},
   description={\textit{Potência}}
}
%%%%%
\newglossaryentry{Pmax}
{
   name={$P_{max}$},
   description={\textit{Potência Máxima}}
}
%%%%%
\newglossaryentry{alpha}
{
   name={$\alpha$},
   description={\textit{Ângulo de Disparo}}
}
%%%%%
\newglossaryentry{pi}
{
   name={$\pi$},
   description={\textit{Constante Pi(\(\approx 3.1416\))}}
}
